\section{Measurement-description overview}
A reusable description specifies one quantum instrument whose behaviour
remains close to the supplied target through every allowed interaction.
This problem depends not only on how many real parameters the target
has, but on how repeated use resolves its different directions.
We study the transition from noisy to extremal readout, together with
singular conditional preparations and the geometry of reusable codes.

For two memoryless instruments of the same interface, define
\[
 d_N(\Phi,\Psi)=\sup_{\mathsf T}
       \|\mathsf T[\Phi]-\mathsf T[\Psi]\|_1.
\]
The same causal tester $\mathsf T$ is used in both hypotheses. It may
have initial reference entanglement, arbitrary finite quantum memory,
feedback and a public stopping time bounded by $N$. The output includes
its retained classical and quantum systems. Thus the distance is
unhalved and has maximum two; event probabilities differ by at most
$d_N/2$. Each device invocation uses the same memoryless instrument.
Covering centres are legal instruments, not arbitrary arrays.

The unbiased-slice result concerns the ordered binary measurements
$E_x^\pm=(I\pm x\cdot\sigma)/2$, $|x|\le1$.
Theorem~\ref{thm:ballmetric74} supplies one finite-use metric comparison
through both the erased and projective boundaries. It treats visibility
as part of the target, rather than a public constant. The proof combines
an exact radial binomial reduction with the retained same-visibility
angular theorem. A monotonicity argument prevents cancellation between
radial and angular perturbations; a simple triangle estimate alone
would not establish the lower bound.

Theorem~\ref{thm:weightedcover74} turns this comparison into a geometric
criterion on any compact Ahlfors-regular identifiable subset of the ball.
Its weighted volume counts the number of local resolution cells. If
the measure of the depth-$t$ boundary layer has order $t^\alpha$,
Corollary~\ref{cor:contact74} distinguishes $\alpha>k/2$, $\alpha=k/2$
and $\alpha<k/2$. Thus the horizon exponent depends on boundary contact,
not just parameter dimension. For the disk the critical equality gives
an additional $\log N$ in the codebook size; the ball has boundary-dominated
order $N^2$. The rational construction in
Theorem~\ref{thm:jointcodec74} realizes both orders with a single charged
index and no irrational normalization of the supplied vector.
These are small-error covering theorems with arbitrary legal lower
centres and family-valued upper centres. The biased-body extension is given in
Theorems~\ref{thm:biasedmetric75} and~\ref{thm:biasedcover75};
the regular-subfamily criterion here retains its unbiased-slice hypotheses.

Single-shot discrimination of white-noise mixtures, including distinct
noise levels, is explicitly treated by Sedl\'ak--Ziman
\cite[Section VI, Example 4]{SZ74}. The optimized single-shot projective
measurement distance is identified by Pucha\l a et al.
\cite[Theorem 1]{PPKK74}. We do not claim those reductions, classical
simulation, or the metrological noise crossover as new. The new
conclusions are the finite-use joint metric, weighted family-covering
criterion, contact law and exact rational description.

We next recall the earlier fixed-visibility and instrument results,
whose complete proofs are retained. The fixed-visibility family consists of the labelled qubit readouts
\[
 E_{\theta,\lambda}^{\pm}=\tfrac12\bigl(I\pm\lambda
              (\cos\theta\,\sigma_1+\sin\theta\,\sigma_2)\bigr).
\]
The visibility $\lambda\in[0,1]$ is public and may vary with $N$.
Theorem~\ref{thm:noisycoding73} gives, with absolute constants,
\[
 d_N(\mathcal M_{\theta,\lambda},\mathcal M_{\phi,\lambda})
 \asymp \min\{1,K_{N,\lambda}h(\theta,\phi)\},\qquad
 K_{N,\lambda}=\lambda\sqrt{N\min\{N,(1-\lambda^2)^{-1}\}}.
\]
Consequently the optimal supplied description has length
$\log_2(1+K_{N,\lambda}/\delta)+O(1)$ for
$0<\delta\le2^{-10}$, uniformly in $N,\lambda,\delta$.
At zero visibility the family is exactly a singleton. At unit
visibility its readout resolution is $N^{-1}$. Between these endpoints,
$N(1-\lambda^2)$ determines the transition. In particular, a noise
sequence $1-\lambda_N^2=N^{-\alpha+o(1)}$ approaching zero gives
horizon coefficient $(1+\min\{\alpha,1\})/2$ at fixed small error.
This is a statement about the operational covering metric, not a
universal parameter-count rule at a boundary.

The upper estimate represents each pair by common projective endpoints
with two classical programme laws and retains their exact finite-use
fidelity. It controls reference assistance and all adaptive testers.
The lower estimate uses blocks of size comparable to
$\min\{N,(1-\lambda^2)^{-1}\}$, or shorter blocks when angular
separation requires them. An explicit Bernoulli product estimate then
gives a finite-distance witness. Noise-dependent loss of Heisenberg
scaling and classical channel simulation are established metrological
methods \cite{DKG67}; their use here is credited. A Fisher-information
bound alone would not imply the displayed global pairwise metric law,
arbitrary-centre covering, or its rational realization.

Two rational semicircle charts yield exact positive effects and a
finite grid decision using only integer arithmetic. The lower cover
allows arbitrary legal memoryless centres. Appending labelled
conditional states of ranks at most $r_+,r_-$ gives
Corollary~\ref{cor:noisypreparation73}: the length acquires the term
$(V/2)\log_2N+V\log_2(1/\delta)$, where
$V=\sum_{y=\pm}[r_y(2n-r_y)-1]$.
The maximally mixed input isolates those state coordinates uniformly
in visibility, while discarding their quantum outputs isolates readout.
The rational upper code preserves each conditional rank. No arbitrary
centre retraction or full large-error theorem is inferred for this
fixed-visibility family.

The previously established observable projective family is
\[
 \mathcal F_{P,\sigma,(x,y)}(X)=\tr(P_xX)\sigma_{xy}.
\]
Here the $P_x$ form an ordered rank-one projective measurement, both
labels $(x,y)$ are retained, and each conditional row of positive
blocks is trace normalized. The rank bounds $r_{xy}$ allow changing
supports and zero blocks. Put $b=d^2-d$ and
$V=\sum_x(\sum_y r_{xy}(2n-r_{xy})-1)$.
Theorem~\ref{thm:readoutentropy72} proves
\[
 \log_2\mathcal R_N(\delta)
  =(b+V/2)\log_2N+(b+V)\log_2(1/\delta)+O(1),
 \qquad 0<\delta\le1/32.
\]
The remainder is independent of $N\ge1$ and $\delta$, at fixed
public dimensions and ranks. The lower bound allows arbitrary legal
memoryless centres; the upper code returns rational family members
without increasing conditional ranks. This is an observable-readout
theorem, not a dimension count for all entanglement-breaking channels.
Corollary~\ref{cor:observablereadout72} removes the explicit $x$ flag
when fixed orthogonal output supports recover it. Equal overlapping
unlabelled rows, by contrast, can erase every basis parameter.

The basis scale follows from the established multiple-query theorem
for von Neumann measurements \cite{PPKKO72}, not from the intrinsic
dimension of the flag manifold alone. Readout differences remain
visible when the conditional outputs are discarded. Conditional-state
differences at a common readout are visible by a basis-state input.
These two pair-dependent witnesses justify a product packing at scales
$\delta/N$ and $\delta/\sqrt N$. For the upper code, partial-pivot
elimination selects bounded rational coordinates directly from the
ordered projections. Normalized algebraic eigenvectors and exhaustive
unitary-net search are unnecessary. Positive row-factor codes and a
common programme comparison then control all adaptive testers.

The inherited Theorem~\ref{thm:seizablecover72}, identifies the
covering problem exactly when one common tester extracts the programme
from one use throughout its ambient state domain. State/channel
comparisons under environment parametrization and seizing are classical
\cite{DW72,WBHK72,WW72}. The additional centre argument retracts every
legal memoryless centre through the same seizer and processor. Thus
finite-use family covering numbers equal their state counterparts,
not merely a binary asymptotic exponent. The inherited rank-state law
then gives the full range $0<\delta\le\delta_*<2$ for fixed cap;
Corollary~\ref{cor:paulicoding72} illustrates it with Pauli channels.
The retraction may increase Choi rank. Rank and zero preservation
belong to the upper target encoder, not to arbitrary-centre retraction.

The preceding results are placed alongside their necessary preliminary
geometry. The fixed-readout input-dependent theorem
\ref{thm:cqentropy71} has coefficient $V/2$ in $\log N$ and the full
fixed submaximal-error range. The common-estimator theorem
\ref{thm:fullpacking71} explains that range by disjoint-event
multiplicity, rather than pairwise separation beyond half the diameter.
The coherent/preparation tensor theorem
\ref{thm:coherententropy70} has projective-unitary dimension $u=d^2-1$
and coefficient $u+v/2$; its small-error range is unchanged.
These earlier theorems, the intrinsic Choi codec and the singular
preparation code are retained with complete proofs. In particular the
new readout quotient removes all column phases, not merely one global
phase. Revision genealogy is recorded outside the article.

The encoder receives matrix data. Fixed dimensions, rank bounds,
horizon, tolerance, error cap and, in the noisy-readout theorem,
visibility are public. Finite chart headers,
pivot masks and coordinate digits are charged in the reusable payload;
expanded matrices, encoder workspace and quantum programme hardware
are different resources. Exact rational implementation applies to
rational targets; the real-target theorem constructs rational centres
in exact real arithmetic. No theorem here learns an unknown device
from queries or implements its action with a finite classical message.
